% =========================================================================
% PLATFORM TAXONOMY ROUTING (patho.dentnet.space)
% =========================================================================
% DOC_ID: developmental_lesions
% DOC_SLUG: developmental_lesions
% DOC_CATEGORY: cat_developmental
% DOC_TITLE: الگوریتم جامع: ضایعات تکاملی و سلول غول‌آسای فکین
% DOC_ALIASES: developmental lesions, ضایعات تکاملی, جاینت سل, CGCG, کیست آنوریسمال, ABC, تومور براون, شروبیسم, SH3BP2
% =========================================================================
% === CLINICAL EVIDENCE LEDGER & PMIDS ===
% 1. WHO Classification of Head and Neck Tumours (5th Ed, 2024) - Giant Cell Lesions and Bone Cysts.
% 2. Neville BW, Damm DD, Allen CM, Chi AC. Oral and Maxillofacial Pathology (5th Ed).
% 3. PMID: 42361776 - Guimarães LM et al. (Semin Diagn Pathol 2026) Molecular Pathology of Distinctive Giant Cell Lesions.
% 4. PMID: 41456990 - Leonardi N et al. (Oral Surg Oral Med 2026) Bone Regeneration in Expansile Jaw Lesions.
% 5. PMID: 35476304 - Oginni FO et al. (Oral Maxillofac Surg 2023) Targeted Biological Therapies in Jaw Giant Cell Lesions.
% 6. PMID: 38814761 - Denosumab vs Intralesional Triamcinolone in Aggressive CGCG Meta-Analysis.
% =========================================================================

\documentclass[a4paper]{article}
\usepackage[margin=1.15cm,footskip=0pt,headsep=0pt]{geometry}
\usepackage{amsmath,amssymb}
\usepackage{xcolor}
\usepackage{tikz}
\usetikzlibrary{arrows.meta,positioning,shapes.geometric}
\usepackage{hyperref}
\usepackage{fontspec}

\makeatletter
\newcommand{\MathCollectFalse}{}
\newcommand{\MathCollectTrue}{}
\let\@kernel@tabular@init\relax
\makeatother

\usepackage{xepersian}

\settextfont[
  Path = Peyda/,
  BoldFont = Peyda-Bold.ttf,
  ItalicFont = Peyda-Medium.ttf,
  BoldItalicFont = Peyda-Bold.ttf
]{Peyda-Regular.ttf}

\setlatintextfont{Times New Roman}

\IfFileExists{Peyda/Symbola.ttf}{%
  \newfontfamily\symb[Path=Peyda/]{Symbola.ttf}%
}{%
  \IfFileExists{editor/Peyda/Symbola.ttf}{%
    \newfontfamily\symb[Path=editor/Peyda/]{Symbola.ttf}%
  }{%
    \newfontfamily\symb{Symbola}%
  }%
}
\newcommand{\emoji}[1]{{\symb #1}}

\linespread{1.10}
\pagestyle{empty}
\thispagestyle{empty}

\raggedbottom
\setlength{\parindent}{0pt}
\setlength{\parskip}{0pt}

\definecolor{hlgray}{HTML}{E2E8F0}
\definecolor{hlyellow}{HTML}{FEF08A}
\definecolor{hlred}{HTML}{FEE2E2}
\definecolor{darkred}{HTML}{991B1B}
\definecolor{navy}{HTML}{0A2540}
\definecolor{teal}{HTML}{0D9488}

\begin{document}

\begin{center}
{\Large\bfseries فلوچارت تصمیم‌گیری بالینی: ضایعات تکاملی و سلول غول‌آسای فکین}\\[2pt]
{\normalsize\bfseries\color{navy} \lr{Developmental and Giant Cell Lesions of the Jaws (Differential Triage)}}\\[3pt]
{\small\itshape تریاژ اندوکرین (\lr{PTH})، سنجش ژن‌های \lr{SH3BP2} و \lr{USP6}، تزریق استروئید داخل ضایعه و درمان‌های هدفمند بیولوژیک}
\end{center}
\vspace{-4pt}
\hrule height 0.8pt
\vspace{6pt}

\begin{center}
\begin{latin}
\linespread{0.95}\selectfont
\begin{tikzpicture}[
    every node/.style={font=\sffamily\footnotesize, align=center},
    stepbox/.style={rectangle, draw=black, fill=white, line width=0.7pt, inner sep=4.5pt},
    decision/.style={diamond, aspect=2.4, draw=black, fill=white, line width=0.7pt, inner sep=2pt, font=\sffamily\scriptsize\bfseries},
    arrow/.style={-{Stealth[scale=0.9]}, thick, black}
]

% 1. Initial Presentation
\node (n1) [stepbox, line width=1.0pt, text width=12.2cm] at (0, 0.0) {
    \rl{\textbf{تظاهر بالینی و رادیوگرافیک: اتساع استخوان فک با رادیولوسنسی اکسپنسیل تک یا چندحفره‌ای}}\\[2pt]
    \scriptsize \rl{اتساع مندیبل یا ماگزیلا $\bullet$ جابجایی یا تحلیل ریشه دندان‌ها $\bullet$ توده با قوام الاستیک یا نمای حباب صابونی}
};

% 2. Mandatory Endocrine Screening
\node (n2) [stepbox, text width=14.2cm] at (0, -1.60) {
    \rl{\colorbox{hlyellow}{\strut\textbf{مرحله ۱ الزامی: تریاژ آزمایشگاهی اندوکرین (سنجش \lr{iPTH} و کلسیم سرم) و آسپیراسیون سوزنی}}}\\[2pt]
    \scriptsize \rl{رد قطعی تومور براون هیپرپاراتیروئیدیسم $\bullet$ ارزیابی آسپیراسیون مایع خونی (مشکوک به کیست آنوریسمال یا مالفورماسیون عروقی)}
};

% 3. Hyperparathyroidism Decision
\node (d1) [decision] at (0, -3.30) {\rl{افزایش}\\[-1pt]\scriptsize \lr{PTH / Ca?}};
\node (r_benign) [stepbox, draw=teal, line width=0.9pt, text width=3.8cm] at (5.8, -3.30) {
    \rl{\textbf{تومور براون (\lr{HPT}) \emoji{⚖}}}\\[2pt]
    \scriptsize \rl{افزایش پاراتورمون و کلسیم سرم؛}\\\rl{ارجاع فوری به جراحی اندوکرین؛}\\\rl{ترمیم خودبه‌خودی استخوان فک}
};

% 4. Mandatory Biopsy & Molecular Assessment
\node (n3) [stepbox, line width=1.0pt, text width=15.8cm] at (0, -5.00) {
    \rl{\colorbox{hlyellow}{\strut\textbf{بیوپسی برشی الزامی (\lr{Incisional Biopsy}) و سنجش ژنتیک \lr{USP6} و \lr{SH3BP2} و مپینگ سه‌بعدی \emoji{🔬}}}}\\[2pt]
    \scriptsize \rl{تفکیک سلول‌های غول‌آسا چند‌هسته‌ای در استرومای فیبروبلاستی $\bullet$ بررسی بازآرایی \lr{USP6} در \lr{ABC} و موتاسیون \lr{SH3BP2} در شروبیسم \lr{[\textbf{PMID: 42361776}]}}
};

% 5. Subtype Decision
\node (d2) [decision] at (0, -6.60) {\rl{فنوتیپ و پاتولوژی}\\[-1pt]\scriptsize \lr{Clinical Phenotype?}};

% 6. 4 Main Columns
\node (h1) [stepbox, text width=3.6cm] at (-6.15, -9.00) {
    \rl{\textbf{\lr{CGCG} غیرتهاجمی}}\\[2pt]
    \scriptsize \lr{Non-Aggressive CGCG}\\\rl{رشد آهسته، کوچکتر از ۲ سانتی‌متر}\\\rl{کورتکس سالم؛ بدون تحلیل ریشه}\\\rl{عود با کورتاژ: \colorbox{hlgray}{\strut\textbf{\rl{کمتر از ۱۰٪}}}}
};
\node (h2) [stepbox, draw=darkred, line width=1.2pt, text width=3.6cm] at (-2.05, -9.00) {
    \rl{\textbf{\lr{CGCG} تهاجمی \emoji{⚡}}}\\[2pt]
    \scriptsize \lr{Aggressive CGCG}\\\rl{رشد سریع، درد و پرفوراسیون کورتکس}\\\rl{تحلیل ریشه و جابجایی دندان‌ها}\\\rl{عود با کورتاژ ساده: \colorbox{hlred}{\strut\color{darkred}\textbf{۴۰٪ تا ۷۰٪}}}
};
\node (h3) [stepbox, line width=0.9pt, text width=3.6cm] at (2.05, -9.00) {
    \rl{\textbf{کیست آنوریسمال (\lr{ABC})}}\\[2pt]
    \scriptsize \lr{Aneurysmal Bone Cyst}\\\rl{فضاهای عروقی پر از خون فاقد اندوتلیوم}\\\rl{بازآرایی ژن \lr{USP6} نئوپلاستیک}\\\rl{عود با کورتاژ ساده: \textbf{۲۰٪ تا ۳۰٪}}
};
\node (h4) [stepbox, line width=0.9pt, text width=3.6cm] at (6.15, -9.00) {
    \rl{\textbf{شروبیسم فامیلیال}}\\[2pt]
    \scriptsize \lr{Cherubism (SH3BP2)}\\\rl{اتساع دوطرفه قرینه زاویه فکین در اطفال}\\\rl{چهره فرشته‌گونه و نمای چشم به آسمان}\\\rl{پسرفت خودبه‌خودی: \colorbox{hlgray}{\strut\textbf{\rl{پس از بلوغ}}}}
};

% 7. Primary Interventions
\node (i1) [stepbox, text width=3.6cm] at (-6.15, -11.60) {
    \rl{\textbf{کورتاژ + استکتومی}}\\[2pt]
    \scriptsize \rl{کورتاژ کامل بافت ضایعه}\\\rl{استکتومی محیطی ۱ تا ۲ میلی‌متر}\\\rl{حفظ دندان‌ها و استخوان فک}
};
\node (i2) [stepbox, draw=darkred, line width=1.2pt, text width=3.6cm] at (-2.05, -11.60) {
    \rl{\textbf{تزریق استروئید یا دنوزوماب}}\\[2pt]
    \scriptsize \textbf{\rl{تریامسینولون هفتگی داخل ضایعه}}\\\rl{به مدت ۶ هفته، یا رزکسیون مارجینال}\\\rl{دنوزوماب در موارد شدید \lr{[\textbf{PMID: 35476304}]}}
};
\node (i3) [stepbox, line width=0.9pt, text width=3.6cm] at (2.05, -11.60) {
    \rl{\textbf{کورتاژ + کرایوتراپی}}\\[2pt]
    \scriptsize \rl{تخلیه کامل حفرات عروقی}\\\rl{کاربرد ادجوانت کرایوتراپی یا فنل}\\\rl{اسکلروتراپی با داکسی‌سایکلین در موارد وسیع}
};
\node (i4) [stepbox, line width=0.9pt, text width=3.6cm] at (6.15, -11.60) {
    \rl{\textbf{مشاهده و مدیریت انتظاری}}\\[2pt]
    \scriptsize \textbf{\rl{پرهیز از جراحی در فاز فعال رشد}}\\\rl{کانتورینگ صورت فقط در صورت اختلال شدید}\\\rl{به تعویق انداختن مداخله تا پایان بلوغ}
};

% 8. Recall Schedules
\node (s1) [stepbox, text width=3.6cm] at (-6.15, -14.20) {
    \rl{\textbf{زمان‌بندی پیگیری:}}\\[2pt]
    \colorbox{hlgray}{\strut\textbf{\rl{هر ۶ ماه تا ۲ سال}}}\\[2pt]
    \scriptsize \rl{تصویربرداری پانورامیک دوره‌ای؛}\\\rl{تایید استخوان‌سازی بستر}
};
\node (s2) [stepbox, draw=darkred, line width=1.2pt, text width=3.6cm] at (-2.05, -14.20) {
    \rl{\textbf{زمان‌بندی پیگیری:}}\\[2pt]
    \colorbox{hlred}{\strut\color{darkred}\textbf{\rl{هر ۳ تا ۴ ماه (۲ سال)}}}\\[2pt]
    \scriptsize \rl{پایش دقیق با \lr{CBCT}؛}\\\rl{بررسی پرفوراسیون و عود زودهنگام}
};
\node (s3) [stepbox, line width=0.9pt, text width=3.6cm] at (2.05, -14.20) {
    \rl{\textbf{زمان‌بندی پیگیری:}}\\[2pt]
    \colorbox{hlgray}{\strut\textbf{\rl{هر ۶ ماه تا ۳ سال}}}\\[2pt]
    \scriptsize \rl{پایش پرشدن لومن عروقی با استخوان؛}\\\rl{تصویربرداری دوره‌ای فک}
};
\node (s4) [stepbox, line width=0.9pt, text width=3.6cm] at (6.15, -14.20) {
    \rl{\textbf{زمان‌بندی پیگیری:}}\\[2pt]
    \colorbox{hlgray}{\strut\textbf{\rl{سالانه تا پایان بلوغ}}}\\[2pt]
    \scriptsize \rl{فتوگرافی سالانه کانتور صورت؛}\\\rl{پایش مسیر رویش دندان‌ها و تنفس}
};

% 9. Trigger Box
\node (trigger) [stepbox, line width=0.9pt, text width=16.4cm, align=right] at (0, -17.40) {
    \centerline{\rl{\textbf{هشدارهای بالینی و اندیکاسیون‌های ارجاع فوری یا مداخله جراحی مجدد (\lr{Clinical Red Flags}):}}}\par\vspace{3pt}
    \scriptsize \rl{\colorbox{hlyellow}{\strut\textbf{مداخله فوری تشخیصی-درمانی در:}} (۱) رشد سریع غیرعادی یا پرفوراسیون کورتکس باکال/لینگوال $\bullet$ (۲) خونریزی ضربان‌دار حاد حین بیوپسی یا آسپیراسیون سوزنی (خطر مالفورماسیون عروقی با فلوی بالا) $\bullet$ (۳) اختلال تنفسی یا دیسفاژی ناشی از اتساع توده شروبیسم به راه هوایی یا کام $\bullet$ (۴) بروز بحران هیپرکلسمی بازگشتی (\lr{Rebound Hypercalcemia}) پس از قطع دنوزوماب $\bullet$ (۵) شکستگی پاتولوژیک فک ناشی از استئولیز وسیع.}
};

% 10. Refractory Box
\node (refract) [stepbox, text width=16.4cm, align=right] at (0, -20.00) {
    \centerline{\rl{\textbf{پروتکل موارد مقاوم، درمان‌های نوین هدفمند و پایش تومور جاینت‌سل (\lr{Malignancy Surveillance}):}}}\par\vspace{3pt}
    \scriptsize \rl{\textbf{سطح ۱ (مهار هدفمند \lr{RANKL}):} کاربرد دنوزوماب زیرجلدی (\lr{120\,mg}) در روزهای ۱، ۸، ۱۵، ۲۹ و سپس هر ۴ هفته در ضایعات راجعه و وسیع با پایش دقیق کلسیم و تیپرینگ تدریجی جهت پیشگیری از هیپرکلسمی بازگشتی.}\par\vspace{2.5pt}
    \scriptsize \rl{\textbf{سطح ۲ (مهار آنزیمی مولکولی):} مهارکننده مسیر \lr{MAPK} (\lr{Trametinib}) در بیماران مبتلا به جهش‌های \lr{TRPV4/KRAS} با ضایعات اگرسیو مقاوم به استروئید.}\par\vspace{2.5pt}
    \scriptsize \rl{\textbf{سطح ۳ (پایش بدخیمی):} بیوپسی برشی مجدد در صورت تغییر رفتار بیولوژیک جهت رد تومور ژانت سل بدخیم فک (\lr{Malignant Giant Cell Tumor}) یا استئوسارکوم غنی از ژانت سل.}
};

% Connections
\draw [arrow] (n1) -- (n2);
\draw [arrow] (n2) -- (d1);
\draw [arrow] (d1.east) -- node[above, font=\scriptsize\bfseries] {\rl{خاموش / خوش‌خیم}} (r_benign.west);
\draw [arrow] (d1.south) -- node[right=5pt, font=\scriptsize\bfseries] {\rl{فعال / مشکوک}} (n3.north);
\draw [arrow] (n3.south) -- (d2.north);

% Bus-Bar from Decision 2 to 4 Subtype columns
\draw [thick, black] (d2.south) -- (0, -7.50);
\draw [thick, black] (-6.15, -7.50) -- (6.15, -7.50);
\draw [arrow] (-6.15, -7.50) -- (h1.north);
\draw [arrow] (-2.05, -7.50) -- (h2.north);
\draw [arrow] (2.05, -7.50) -- (h3.north);
\draw [arrow] (6.15, -7.50) -- (h4.north);

% Column Downward Arrows
\draw [arrow] (h1.south) -- (i1.north);
\draw [arrow] (h2.south) -- (i2.north);
\draw [arrow] (h3.south) -- (i3.north);
\draw [arrow] (h4.south) -- (i4.north);

\draw [arrow] (i1.south) -- (s1.north);
\draw [arrow] (i2.south) -- (s2.north);
\draw [arrow] (i3.south) -- (s3.north);
\draw [arrow] (i4.south) -- (s4.north);

% Collector Bus into Trigger Box
\draw [thick, black] (s1.south) -- (-6.15, -15.80);
\draw [thick, black] (s2.south) -- (-2.05, -15.80);
\draw [thick, black] (s3.south) -- (2.05, -15.80);
\draw [thick, black] (s4.south) -- (6.15, -15.80);
\draw [thick, black] (-6.15, -15.80) -- (6.15, -15.80);
\draw [arrow] (0, -15.80) -- (trigger.north);

% Final link to Refractory Box
\draw [arrow] (trigger.south) -- (refract.north);
\end{tikzpicture}
\end{latin}
\end{center}

\end{document}
